\documentclass[11pt]{article}
\usepackage[a4paper,margin=1in]{geometry}
\usepackage[T1]{fontenc}
\usepackage{lmodern}
\usepackage{amsmath,amssymb,amsthm,booktabs,enumitem,mathtools,microtype}
\usepackage[numbers,sort&compress]{natbib}
\usepackage{xurl}
\usepackage[hidelinks]{hyperref}
\setlength{\emergencystretch}{3em}

\newtheorem{theorem}{Theorem}[section]
\newtheorem{proposition}[theorem]{Proposition}
\newtheorem{lemma}[theorem]{Lemma}
\newtheorem{corollary}[theorem]{Corollary}
\theoremstyle{definition}
\newtheorem{definition}[theorem]{Definition}
\theoremstyle{remark}
\newtheorem{remark}[theorem]{Remark}

\newcommand{\mob}{\mu}
\newcommand{\ind}{\mathbf 1}
\newcommand{\T}{\{-1,0,+1\}}
\newcommand{\Z}{\mathbb Z}
\newcommand{\Target}{\operatorname{Target}}
\newcommand{\Source}{\operatorname{Source}}
\newcommand{\rad}{\operatorname{rad}}

\hypersetup{
  pdftitle={Euler Run Spectrum for Fixed-Lag Centered Mobius Capacity},
  pdfauthor={RH research program},
  pdfsubject={Fixed-lag centered Mobius capacity, Euler-run endpoints, and finite-clock nonattainment},
  pdfkeywords={Mobius function, terminal logarithmic average, fixed lag,
  max-plus optimization, squarefree runs, Euler product}
}

\title{Euler--Run Spectrum for Fixed-Lag Centered M\"obius Capacity}
\author{RH research program}
\date{August 12, 2026}

\begin{document}
\maketitle

\begin{abstract}
Fix a lag \(h\geq1\).  A centered phase table reads
\(\mob_0(n-h),\mob(n),\mob(n+h)\), and universal safety forbids two
positive outputs at separation \(2h\) for every ternary word.  The
complete three-shift terminal-log law converts every fixed-table limit,
for every terminal clock, into collision-aware squarefree densities.
Positive projection and relation saturation then give an exact tropical
trace on all eight subsets of \(\{-1,0,+1\}\) for every finite clock.
The familiar four-state compression is proved throughout the
non-self-loop regime \(q\nmid2h\); when \(q\mid2h\), the full eight states
remain necessary in general.

On square-support clocks, a phasewise identity for each shared ternary
value charges two consecutive transitions by one center weight.  Once the
base support contains
\(p_0(h)=\min\{p\text{ prime}:p\nmid2h\}\), the centered capacity equals a
raw step-\(2h\) path-MWIS normalized by the squarefree density.  Exact
bracketed-run densities yield the Euler--run endpoint
\[
 B_\infty(h)=\frac{3}{\pi^2}
 +\frac12\sum_{\substack{1\leq\ell<p_0(h)^2\\ \ell\ \mathrm{odd}}}
 R_{\ell,h}.
\]
For every fixed \(h\), this is the supremum over finite clocks, and no
finite clock attains it.  Fresh-prime lifts can plateau: strict gain is
equivalent to the presence of an even positive run, while a CRT
construction guarantees eventual strict gain.  Across fixed lags,
\(\inf_{h\geq1}B_\infty(h)=3/\pi^2\), and the infimum is not attained.
No monotonicity, supremum, or maximum in \(h\) is claimed.
\end{abstract}

\noindent\textbf{Keywords:} M\"obius function; terminal logarithmic
average; centered local rule; fixed lag; tropical trace; squarefree run;
nonattained supremum.

\section{Definitions and main results}

Write \(\mob_0(k)=\mob(k)\) for integers \(k\geq1\), and
\(\mob_0(k)=0\) for \(k\leq0\).  A \emph{terminal clock} is a function
\begin{equation}
 1\leq\omega(X)\leq X,
 \qquad \omega(X)\longrightarrow\infty.
 \label{eq:clock}
\end{equation}
Fix \(h,q\geq1\), put \(d=2h\), and fix, before \(X\to\infty\), one sign
table for every phase,
\[
 F_r:\T^3\longrightarrow\{-1,+1\},\qquad r\in\Z/q\Z.
\]
Its centered output and scored terminal functional are
\begin{align}
 \varepsilon_F(n)
 &=F_{n\bmod q}\bigl(\mob_0(n-h),\mob(n),\mob(n+h)\bigr),
 \label{eq:output}\\
 L_{h,q,X}(F)
 &=\frac1{\log\omega(X)}
 \sum_{X/\omega(X)<n\leq X}\frac{\mob(n)\varepsilon_F(n)}n.
 \label{eq:functional}
\end{align}
Thus \(\varepsilon_F(n)\) is the output, while
\(\mob(n)\varepsilon_F(n)\) is the score.  The family is
\emph{universally distance-\(d\) safe} when
\begin{equation}
 \neg\bigl(F_r(a,b,c)=+1\ \text{and}\
 F_{r+d}(c,e,f)=+1\bigr)
 \label{eq:safety}
\end{equation}
for every phase \(r\) and every \(a,b,c,e,f\in\T\).  The shared letter is
\(c=\mob(n+h)=\mob((n+d)-h)\).

The density notation must retain collisions among the three shifts.  Let
\(\mathcal I=\{L,C,R\}\) and
\begin{equation}
 a_L=h,\qquad a_C=0,\qquad a_R=-h.
 \label{eq:coordinate-shifts}
\end{equation}
For \(S\subseteq\mathcal I\), first deduplicate modulo \(p^2\), then set
\begin{equation}
 B_{p,S}=\{a_i\bmod p^2:i\in S\},\qquad
 \nu_{p,S}=|B_{p,S}|,
 \qquad
 \tau_{p,S}(r)=\#\{b\in B_{p,S}:b\equiv r\pmod p\}.
 \label{eq:local-data}
\end{equation}
Define
\begin{align}
 \Theta_{h,q,r}(S)
 &=\frac1q
 \prod_{p\nmid q}\left(1-\frac{\nu_{p,S}}{p^2}\right)
 \prod_{p\parallel q}\left(1-\frac{\tau_{p,S}(r)}p\right)
 \prod_{p^2\mid q}
 \ind_{\{r\bmod p^2\notin B_{p,S}\}},
 \label{eq:theta}\\
 \kappa_h(S)
 &=\prod_p\left(1-\frac{\nu_{p,S}}{p^2}\right).
 \label{eq:kappa}
\end{align}
Then \(\Theta_{h,q,r}(\varnothing)=1/q\) and
\begin{equation}
 \sum_{r\bmod q}\Theta_{h,q,r}(S)=\kappa_h(S).
 \label{eq:phase-sum}
\end{equation}
Only in a collision-free \(j\)-site case does this reduce to
\begin{equation}
 K_j=\prod_p\left(1-\frac{j}{p^2}\right),\qquad
 K_0=1,\qquad K_1=\frac6{\pi^2}.
 \label{eq:Kj}
\end{equation}
The phase-resolved exact-support density is
\begin{equation}
 \Pi_{h,q,r}(U)=
 \sum_{W\subseteq\mathcal I\setminus U}(-1)^{|W|}
 \Theta_{h,q,r}(U\cup W).
 \label{eq:Pi}
\end{equation}
It is nonnegative and has mass
\(\sum_{U\subseteq\mathcal I}\Pi_{h,q,r}(U)=1/q\).

For \(x,y\in\T\), write
\begin{equation}
 S(x,y)=\{C\}\cup\bigl(\{L\}:x\ne0\bigr)
 \cup\bigl(\{R\}:y\ne0\bigr)
 \label{eq:Sxy}
\end{equation}
and define the nonnegative cell weight and transition
\begin{align}
 \lambda_{h,q,r}(x,y)
 &=2^{-\ind_{\{x\ne0\}}-\ind_{\{y\ne0\}}}
 \Pi_{h,q,r}\bigl(S(x,y)\bigr),
 \label{eq:lambda}\\
 \mathcal K_{h,q,r}(U,V)
 &=\sum_{\substack{x\notin U\\y\in V}}
 \lambda_{h,q,r}(x,y),
 \qquad U,V\subseteq\T.
 \label{eq:transition}
\end{align}

\begin{theorem}[Fixed-clock terminal law and full tropical formula]
\label{thm:fixed-clock}
For every fixed \(h,q,F\), the limit
\(L_{h,q}(F)=\lim_{X\to\infty}L_{h,q,X}(F)\) exists for every terminal
clock \eqref{eq:clock} and has the same value for all such clocks.  Form
the finite maximum only after these fixed-table limits:
\begin{equation}
 C_h(q)=\max_{F\ {\rm universally\ safe}}|L_{h,q}(F)|.
 \label{eq:capacity}
\end{equation}
If \(\gamma\) runs over the \(\gcd(q,d)\) cycles of
\(r\mapsto r+d\) in \(\Z/q\Z\), then
\begin{equation}
 \boxed{
 C_h(q)=\sum_\gamma
 \max_{\substack{Y_r\subseteq\T\\r\in\gamma\ {\rm cyclic}}}
 \sum_{r\in\gamma}\mathcal K_{h,q,r}(Y_{r-d},Y_r).}
 \label{eq:full-eight}
\end{equation}
Thus the exact formula is a tropical trace on eight states for every
finite \(q\).  Both signs of every nonzero optimum are attained.
\end{theorem}

\begin{theorem}[Four-state scope and self-loop obstruction]
\label{thm:compression}
If \(q\nmid d\), an optimizer in \eqref{eq:full-eight} can be chosen from
the four antipodally symmetric subsets of \(\T\).  If \(q\mid d\), every
cycle is a self-loop and the theorem retains all eight states; no
four-state reduction is asserted.  This exception is strict in general:
for \(h=2,q=4\), the full and four-state-restricted coefficient vectors in
the basis \((K_0,K_1,K_2,K_3)\) are respectively
\begin{equation}
 (0,0,\tfrac12,-\tfrac12)
 \quad\text{and}\quad
 (0,0,1,-2),
 \label{eq:self-loop-obstruction}
\end{equation}
and the full value is strictly larger.
\end{theorem}

For the endpoint, let
\begin{equation}
 p_0(h)=\min\{p\text{ prime}:p\nmid d\}.
 \label{eq:p0}
\end{equation}
For a finite set \(J\subset\Z\), define the absolutely convergent local
product
\begin{equation}
 D_h(J)=\prod_p\left(
 1-\frac{|\{dj\bmod p^2:j\in J\}|}{p^2}\right),
 \label{eq:Dh}
\end{equation}
and, for \(\ell\geq1\),
\begin{align}
 R_{\ell,h}
 &=D_h([0,\ell-1])-D_h(\{-1\}\cup[0,\ell-1])
 \notag\\
 &\quad-D_h([0,\ell])+D_h([-1,\ell]).
 \label{eq:Rlh}
\end{align}
Intervals here denote sets of integers.  The quantity \(R_{\ell,h}\) is
the density of an exact bracketed squarefree run of length \(\ell\) in
step-\(d\) coordinates, hence is nonnegative.  Set
\begin{equation}
 B_\infty(h)=\frac3{\pi^2}
 +\frac12\sum_{\substack{1\leq\ell<p_0(h)^2\\\ell\ \mathrm{odd}}}
 R_{\ell,h}.
 \label{eq:Binfinity}
\end{equation}

\begin{theorem}[Euler--run endpoint and finite nonattainment]
\label{thm:endpoint}
For every fixed \(h\geq1\),
\begin{equation}
 \boxed{\sup_{q<\infty}C_h(q)=B_\infty(h),\qquad
 C_h(q)<B_\infty(h)\quad(q<\infty).}
 \label{eq:all-clock-endpoint}
\end{equation}
The supremum is over finite clocks only and is taken after the fixed-table
terminal limits and finite safe maxima.
\end{theorem}

\begin{corollary}[Lag infimum]
\label{cor:lag-infimum}
Across fixed positive integer lags,
\begin{equation}
 \boxed{\inf_{h\geq1}B_\infty(h)=\frac3{\pi^2}.}
 \label{eq:lag-infimum}
\end{equation}
Every fixed \(h\) satisfies \(B_\infty(h)>3/\pi^2\), so the infimum is
not attained.  No supremum, maximum, or monotonicity assertion in \(h\)
is made.
\end{corollary}

\section{Terminal table bridge and relation saturation}

RH-394 proves the complete terminal-log table law for any three fixed
distinct shifts: Theorem~1.1, equation~(8), Theorem~1.2,
equations~(11)--(14), and Corollary~1.3 (printed/PDF pages 2--3)
\cite{RH394}.  We instantiate it at the ordered shifts
\((a_L,a_C,a_R)=(h,0,-h)\).  They are distinct because \(h\geq1\), and
all of \(h,q,F\) are fixed before the limit.  This is the sole analytic
terminal-clock input in the present paper.

To state the resulting bridge, put
\[
 H_r(x,z,y)=zF_r(x,z,y)
\]
and, for \(U\subseteq\mathcal I\), let
\[
 \overline H_{r,U}=2^{-|U|}
 \sum_{\sigma\in\{-1,+1\}^{U}}H_r(\sigma_U,0_{U^c}).
\]
The cited law, including the boundary convention \(\mob_0\), gives
\begin{equation}
 L_{h,q}(F)=
 \sum_{r\bmod q}\sum_{U\subseteq\mathcal I}
 \Pi_{h,q,r}(U)\overline H_{r,U}.
 \label{eq:table-bridge}
\end{equation}
It also gives \eqref{eq:theta}--\eqref{eq:phase-sum}.  The
deduplication in \eqref{eq:local-data} is essential: if, for example,
\(p^2\mid d\), several coordinate shifts can define the same forbidden
class modulo \(p^2\).  Equation \eqref{eq:table-bridge} is qualitative
and fixed-data; no rate uniform in \(h\) or \(q\) is imported.  Tao and
Tao--Ter\"av\"ainen are inherited two-point and odd-correlation
provenance inside RH-394, not separately strengthened inputs here
\cite{Tao2016LogChowla,TaoTeravainen2019}.

\begin{lemma}[Positive projection]
\label{lem:projection}
Given a safe family \(F\), replace every \(+1\) output at a cell whose
center \(z\) is not \(+1\) by \(-1\), leaving every other cell unchanged.
Call the result \(F^+\).  Then \(F^+\) is safe and
\[
 L_{h,q}(F^+)\geq L_{h,q}(F).
\]
If
\begin{equation}
 A_r=\{(x,y)\in\T^2:F^+_r(x,+1,y)=+1\},
 \label{eq:relation}
\end{equation}
then
\begin{equation}
 L_{h,q}(F^+)=
 \sum_{r\bmod q}\sum_{(x,y)\in A_r}\lambda_{h,q,r}(x,y).
 \label{eq:projected-limit}
\end{equation}
\end{lemma}

\begin{proof}
At \(z=-1\), changing the output from \(+1\) to \(-1\) changes the score
\(zF\) from \(-1\) to \(+1\); at \(z=0\), it does not change the score.
Deleting positive outputs cannot create a forbidden pair, so safety is
preserved.  This establishes the pointwise score comparison before any
limit is taken.

Start now from the all-minus table.  Its score is \(-z\), whose uniform
average on every exact-support sign stratum is zero.  Turning the
center-\(+1\) cell \((x,+1,y)\) from minus to plus changes \(H\) by two.
The density of that specified sign cell is
\[
 2^{-1-\ind_{\{x\ne0\}}-\ind_{\{y\ne0\}}}
 \Pi_{h,q,r}(S(x,y)).
\]
Multiplication by two gives exactly \eqref{eq:lambda}.  Summing the
changed cells proves \eqref{eq:projected-limit}.
\end{proof}

For a relation \(A\subseteq\T^2\), define, before using either symbol,
\[
 \Source(A)=\{x:(x,y)\in A\text{ for some }y\},\qquad
 \Target(A)=\{y:(x,y)\in A\text{ for some }x\}.
\]

\begin{lemma}[Safety and saturation]
\label{lem:saturation}
The projected relation family is safe if and only if
\begin{equation}
 \Target(A_r)\cap\Source(A_{r+d})=\varnothing
 \qquad(r\bmod q).
 \label{eq:relation-safety}
\end{equation}
Put \(Y_r=\Target(A_r)\).  Then
\begin{equation}
 A_r\subseteq(\T\setminus Y_{r-d})\times Y_r.
 \label{eq:relation-containment}
\end{equation}
Replacing every \(A_r\) by the saturated relation
\begin{equation}
 \widetilde A_r=(\T\setminus Y_{r-d})\times Y_r
 \label{eq:saturated-relation}
\end{equation}
preserves safety and weakly increases \eqref{eq:projected-limit}.
Conversely, every subset profile \((Y_r)_{r\bmod q}\) in
\eqref{eq:saturated-relation} is safe.
\end{lemma}

\begin{proof}
An edge \((x,c)\in A_r\) and an edge \((c,y)\in A_{r+d}\) are precisely
two positive outputs allowed by a common shared letter \(c\).  Their
absence for every \(c\in\T\) is \eqref{eq:relation-safety}.  It follows
that every source of \(A_r\) lies outside
\(\Target(A_{r-d})=Y_{r-d}\), while its targets lie in \(Y_r\), proving
\eqref{eq:relation-containment}.  All weights in
\eqref{eq:projected-limit} are nonnegative, so saturation cannot lower
the score.  Finally, an edge of \(\widetilde A_r\) ends in \(Y_r\), while
an edge of \(\widetilde A_{r+d}\) begins outside \(Y_r\); hence the
saturated family is safe.
\end{proof}

\begin{lemma}[Reflection]
\label{lem:reflection}
Define
\[
 F^\rho_r(x,z,y)=F_r(-x,-z,-y).
\]
This map preserves universal safety and satisfies
\begin{equation}
 L_{h,q}(F^\rho)=-L_{h,q}(F).
 \label{eq:reflection-sign}
\end{equation}
\end{lemma}

\begin{proof}
Simultaneous negation is a bijection of the ternary words appearing in
\eqref{eq:safety}, so it preserves safety.  Moreover,
\[
 zF^\rho_r(x,z,y)=-H_r(-x,-z,-y).
\]
Every exact-support stratum in \eqref{eq:table-bridge} is invariant under
simultaneous sign negation.  Its uniform average therefore changes sign,
which proves \eqref{eq:reflection-sign}.  This is a terminal table-law
identity; it does not identify the reflected table with a pointwise
transformation of the observed M\"obius sequence.
\end{proof}

\section{The tropical optimizer and four-state boundary}

\begin{proposition}[Exact eight-state trace]
\label{prop:tropical}
The right side of \eqref{eq:full-eight} is the maximum positive terminal
score over all safe tables.
\end{proposition}

\begin{proof}
By Lemmas~\ref{lem:projection} and~\ref{lem:saturation}, a positive
optimum is attained by a saturated subset profile.  The phase \(r\)
contribution is
\[
 \sum_{\substack{x\notin Y_{r-d}\\y\in Y_r}}
 \lambda_{h,q,r}(x,y)
 =\mathcal K_{h,q,r}(Y_{r-d},Y_r).
\]
Addition by \(d\) decomposes \(\Z/q\Z\) into \(\gcd(q,d)\) cycles, each
of length \(q/\gcd(q,d)\).  Profiles on distinct cycles do not interact,
so their maxima add.  Cyclic closure of each finite dynamic program is
exactly the tropical trace in \eqref{eq:full-eight}.  There are eight
possible subsets of \(\T\), including all singleton states and all
self-loops.
\end{proof}

\begin{proof}[Proof of Theorem~\ref{thm:fixed-clock}]
The existence and clock independence of every fixed-table limit follow
from \eqref{eq:table-bridge}.  Positive projection, saturation, and
Proposition~\ref{prop:tropical} give the displayed positive maximum.
Lemma~\ref{lem:reflection} shows that every score has a safe score of the
opposite sign, hence the positive maximum equals the absolute capacity
and both orientations are attained.  This proves every assertion of the
theorem, including self-loop clocks.
\end{proof}

For an antipodally symmetric set \(Y\subseteq\T\), write
\[
 u(Y)=\left(u_0(Y),u_1(Y)\right)
 =\left(\ind_{\{0\in Y\}},
 \frac{|Y\cap\{-1,+1\}|}{2}\right)\in\{0,1\}^2.
\]
Abbreviate
\[
 \alpha_r=\Pi_{h,q,r}(\{C\}),\quad
 \beta_r=\Pi_{h,q,r}(\{L,C\}),\quad
 \gamma_r=\Pi_{h,q,r}(\{C,R\}),\quad
 \eta_r=\Pi_{h,q,r}(\{L,C,R\}).
\]
Directly summing \eqref{eq:lambda} gives, on the four states,
\begin{align}
 \mathcal K_{h,q,r}(U,V)
 &=\alpha_r(1-u_0(U))u_0(V)
 +\beta_r(1-u_1(U))u_0(V)\notag\\
 &\quad+\gamma_r(1-u_0(U))u_1(V)
 +\eta_r(1-u_1(U))u_1(V).
 \label{eq:four-transition}
\end{align}

\begin{proposition}[Non-self-loop compression]
\label{prop:compression}
If \(q\nmid d\), the eight-state maximum has a four-state optimizer, and
its transitions are given by \eqref{eq:four-transition}.
\end{proposition}

\begin{proof}
Fix the membership of zero in every \(Y_r\), and set
\(k_r=|Y_r\cap\{-1,+1\}|\in\{0,1,2\}\).  The cell weights depend on a
nonzero outer coordinate only through its being nonzero, not through its
sign.  Consequently the two transition terms in which \(Y_r\) occurs
are affine in \(k_r\) once all other memberships are fixed.

The hypothesis \(q\nmid d\) is equivalent to
\(q/\gcd(q,d)>1\).  Thus the occurrence of \(Y_r\) as target in the
\(r\)-term and as excluded source in the \(r+d\)-term belongs to two
distinct terms, with no self-identification.  An affine function at
\(k_r=1\) is no larger than at least one endpoint \(k_r=0\) or \(2\).
Round one phase at a time.  The objective never decreases, and after
finitely many rounds all nonzero memberships are antipodally symmetric.
Equation \eqref{eq:four-transition} follows by splitting the sum into
the four zero/nonzero support patterns.
\end{proof}

\begin{proof}[Proof of Theorem~\ref{thm:compression}]
Proposition~\ref{prop:compression} proves the non-self-loop statement.
If \(q\mid d\), then \(r+d=r\) for every phase.  A self-loop transition
uses the same nonzero count both as an included target and an excluded
source; terms proportional to \(k_r(2-k_r)\) can occur, so the preceding
affine rounding is unavailable.

The obstruction \((h,q)=(2,4)\) can be checked without approximation.
At each even phase every self-loop transition is zero: phase \(0\) has
zero center mass, while at phase \(2\) both outer coordinates are forced
to be zero and cannot lie simultaneously outside and inside the same
state.  At either odd phase, substituting
\eqref{eq:theta}--\eqref{eq:transition} for all eight subsets gives
\begin{center}
\begin{tabular}{ccl}
\toprule
\(Y\) type & number of states & self-loop value\\
\midrule
\(\varnothing,\T\) & \(2\) & \(0\)\\
one nonzero sign, or its complement with \(0\) & \(4\)
  & \((K_2-K_3)/4\)\\
\(\{0\}\) or \(\{-1,+1\}\) & \(2\)
  & \(K_2/2-K_3\)\\
\bottomrule
\end{tabular}
\end{center}
Here
\[
 \frac{K_3}{K_2}
 =\frac12\prod_{p\geq3}\left(1-\frac1{p^2-2}\right).
\]
The product is less than one.  On the other hand,
\[
 \sum_{p\geq3}\frac1{p^2-2}
 <\sum_{\substack{n\geq3\\n\ {\rm odd}}}\frac1{n(n-1)}
 =1-\log2<\frac13,
\]
and \(\prod(1-x_i)\geq1-\sum x_i\), together with
\(\log2>2/3\), gives
\begin{equation}
 \frac13<\frac{K_3}{K_2}<\frac12.
 \label{eq:K3K2-bounds}
\end{equation}
Thus a one-sign state is the full optimum at each odd phase, whereas
the four-state optimum is \(\{0\}\) or \(\{-1,+1\}\).  Adding the two
odd phases yields respectively
\((K_2-K_3)/2\) and \(K_2-2K_3\), which are the coefficient vectors in
\eqref{eq:self-loop-obstruction}; their difference
\((3K_3-K_2)/2\) is positive by \eqref{eq:K3K2-bounds}.
\end{proof}

\section{Square-support marginal charge}

Let \(\mathcal P\) be a nonempty finite set of primes and
\[
 q_{\mathcal P}=\prod_{p\in\mathcal P}p^2.
\]
More generally, call \(Q\) \emph{square-supported} if every prime divisor
of \(Q\) occurs to exponent at least two.  Its positive center phases form
\[
 V_h(Q)=\{r\bmod Q:\Theta_{h,Q,r}(\{C\})>0\}
       =\{r\bmod Q:p^2\nmid r\text{ for every }p\mid Q\}.
\]
Write
\begin{align}
 N_h(Q)&=|V_h(Q)|,\label{eq:NQ}\\
 \alpha_h(Q)&=\max_{\substack{I\subseteq V_h(Q)\\
 I\cap(I+d)=\varnothing}}|I|,\label{eq:alphaQ}\\
 M_h(Q)&=K_1\frac{\alpha_h(Q)}{N_h(Q)}.
 \label{eq:MQ}
\end{align}
The normalization is deliberate: \(\alpha_h(Q)\) is a raw, unweighted
MWIS cardinality, whereas \(M_h(Q)\) is a weighted terminal value.
Neither \eqref{eq:MQ} nor the equality with centered capacity below is
being asserted for clocks having a prime only to the first power.

\begin{lemma}[Clock divisibility and one-site embedding]
\label{lem:clock-divisibility}
If \(q\mid Q\), then \(C_h(q)\leq C_h(Q)\).  On every square-supported
\(Q\), one also has \(M_h(Q)\leq C_h(Q)\).
\end{lemma}

\begin{proof}
Regard a \(q\)-phase table as a \(Q\)-phase table by literal repetition.
It produces the same output word, preserves universal safety, and has the
same fixed-table terminal limit.  The declared clock need not be a
minimal period.  Maximizing after the limits gives the first assertion.

For the second, choose an independent set \(I\) in
\eqref{eq:alphaQ}.  At \(r\in I\), let the projected relation be all of
\(\T^2\); off \(I\), let it be empty.  This family is safe, and summing
\eqref{eq:lambda} over both outer coordinates gives
\(\Theta_{h,Q,r}(\{C\})\).  All positive center phases of a
square-supported clock have the common weight
\begin{equation}
 \delta_Q=\frac1Q\prod_{p\nmid Q}\left(1-\frac1{p^2}\right)
 =\frac{K_1}{N_h(Q)}.
 \label{eq:deltaQ}
\end{equation}
Thus the embedded value is
\(\delta_Q\alpha_h(Q)=M_h(Q)\).
\end{proof}

\begin{lemma}[Per-state shared marginal]
\label{lem:marginal}
Let \(Q\) be square-supported, and suppose that both \(r\) and \(r+d\)
belong to \(V_h(Q)\).  For each \(t\in\T\), define
\[
 m_r(t)=\sum_{x\in\T}\lambda_{h,Q,r}(x,t),\qquad
 m'_{r+d}(t)=\sum_{y\in\T}\lambda_{h,Q,r+d}(t,y).
\]
Then
\begin{align}
 m_r(0)
 &=\Theta_{h,Q,r}(\{C\})-\Theta_{h,Q,r}(\{C,R\})
 \notag\\
 &=\Theta_{h,Q,r+d}(\{C\})
   -\Theta_{h,Q,r+d}(\{L,C\})=m'_{r+d}(0),
 \label{eq:marginal-zero}\\
 m_r(\pm1)
 &=\frac12\Theta_{h,Q,r}(\{C,R\})
 =\frac12\Theta_{h,Q,r+d}(\{L,C\})
 =m'_{r+d}(\pm1).
 \label{eq:marginal-sign}
\end{align}
In particular,
\begin{equation}
 m_r(t)=m'_{r+d}(t)\quad(t\in\T),\qquad
 \sum_{t\in\T}m_r(t)=\delta_Q.
 \label{eq:marginal-total}
\end{equation}
\end{lemma}

\begin{proof}
Summing exact-support cells with right coordinate zero gives the event
that the center is squarefree while the right site is not; this is the
difference in the first line of \eqref{eq:marginal-zero}.  If the right
coordinate is a specified nonzero sign, uniform sign splitting contributes
the factor \(1/2\) in \eqref{eq:marginal-sign}.  The same argument at
phase \(r+d\) uses the left coordinate.

It remains to prove the two phase identities.  Outside the prime support
of \(Q\), the pair of coordinate shifts \(\{C,R\}=\{0,-h\}\) is a
translate of \(\{L,C\}=\{h,0\}\), so the deduplicated local factors are
identical even when the two residues collide.  At a supported prime
square, the first pair requires \(r\not\equiv0,-h\pmod{p^2}\), and the
second requires \(r+d\not\equiv0,h\pmod{p^2}\).  Since the endpoint center
phases \(r\) and \(r+d\) are positive, in both cases the only remaining
condition is \(r+h\not\equiv0\pmod{p^2}\).  Hence
\[
 \Theta_{h,Q,r}(\{C,R\})
 =\Theta_{h,Q,r+d}(\{L,C\}).
\]
Both singleton center densities equal \(\delta_Q\) by
\eqref{eq:deltaQ}.  This proves \eqref{eq:marginal-zero} and
\eqref{eq:marginal-sign}; adding the three states proves
\eqref{eq:marginal-total}.  The \(t=0\) difference is part of the
identity, not an omitted zero-weight case.
\end{proof}

\begin{lemma}[Pair and path charge]
\label{lem:path-charge}
Under the hypotheses of Lemma~\ref{lem:marginal}, for all
\(U,V,W\subseteq\T\),
\begin{equation}
 \mathcal K_{h,Q,r}(U,V)
 +\mathcal K_{h,Q,r+d}(V,W)\leq\delta_Q.
 \label{eq:pair-charge}
\end{equation}
Consequently a positive step-\(d\) path of length \(L\) contributes at
most \(\lceil L/2\rceil\delta_Q\) to
\eqref{eq:full-eight}.
\end{lemma}

\begin{proof}
Dropping the outer-coordinate restrictions in the two transitions and
splitting at the shared state gives
\begin{align*}
 \mathcal K_{h,Q,r}(U,V)
 &\leq\sum_{t\in V}m_r(t),\\
 \mathcal K_{h,Q,r+d}(V,W)
 &\leq\sum_{t\notin V}m'_{r+d}(t).
\end{align*}
Their sum is \(\delta_Q\) by \eqref{eq:marginal-total}.  Pair successive
sites of a path, leaving one unpaired site when \(L\) is odd; a single
transition is at most its full center mass \(\delta_Q\).  This gives the
path charge.
\end{proof}

\begin{proposition}[Square-supported centered saturation]
\label{prop:square-saturation}
If \(Q\) is square-supported and \(p_0(h)\mid Q\), then
\begin{equation}
 \boxed{C_h(Q)=M_h(Q)
 =K_1\frac{\alpha_h(Q)}{N_h(Q)}.}
 \label{eq:square-saturation}
\end{equation}
\end{proposition}

\begin{proof}
Because \(p_0\nmid d\), a step-\(d\) orbit modulo \(p_0^2\) visits every
residue class.  The factor \(p_0^2\mid Q\) therefore forces at least one
zero center phase on every step-\(d\) cycle.  Deleting the zero phases
splits the positive graph into finite paths; there are no all-positive
cycles.  Lemma~\ref{lem:path-charge} bounds the contribution of each path
by its ordinary path-MWIS cardinality times \(\delta_Q\).  Summing paths
gives
\[
 C_h(Q)\leq\delta_Q\alpha_h(Q)=M_h(Q).
\]
The reverse inequality is the one-site embedding from
Lemma~\ref{lem:clock-divisibility}.
\end{proof}

\begin{proposition}[Same-support cover]
\label{prop:same-support}
Suppose \(p_0(h)\in\mathcal P\),
\(q_{\mathcal P}\mid Q\), and
\(\rad(Q)=\rad(q_{\mathcal P})\).  Put
\(R=Q/q_{\mathcal P}\).  No coprimality condition on \(R\) and \(d\) is
required.  Then
\begin{align}
 N_h(Q)&=R\,N_h(q_{\mathcal P}),&
 \alpha_h(Q)&=R\,\alpha_h(q_{\mathcal P}),\label{eq:cover-counts}\\
 C_h(Q)&=M_h(Q)=M_h(q_{\mathcal P})=C_h(q_{\mathcal P}).
 \label{eq:cover-capacity}
\end{align}
\end{proposition}

\begin{proof}
Reduction modulo \(q_{\mathcal P}\) makes the positive support word on
\(\Z/Q\Z\) an \(R\)-fold cover of the corresponding word on
\(\Z/q_{\mathcal P}\Z\), proving the first count in
\eqref{eq:cover-counts}.  By Proposition~\ref{prop:square-saturation},
the base positive graph is a disjoint union of paths.  The inverse image
of a path under any graph cover is \(R\) disjoint copies of that path:
choose a lift of one endpoint, then lift its unique successive edges.
This argument is independent of \(\gcd(R,d)\).  Hence the path-MWIS
cardinality scales by \(R\), proving the second count.  The ratio in
\eqref{eq:MQ} is unchanged, and
Proposition~\ref{prop:square-saturation} at both clocks proves
\eqref{eq:cover-capacity}.
\end{proof}

\begin{remark}[Why the \(p_0\) qualification is necessary]
\label{rem:p0-counterexample}
For \(h=6\), one has \(d=12\) and \(p_0=5\).  The same-support lift
\(36\mid72\) occurs before the support contains \(5^2\):
\[
 (\alpha_6(36),N_6(36))=(9,24),\qquad
 (\alpha_6(72),N_6(72))=(24,48).
\]
The raw ratios are \(3/8\) and \(1/2\), so unconditional
same-support scaling is false.  Once the base includes \(p_0^2\),
\[
 (\alpha_6(900),N_6(900))=(291,576),\qquad
 (\alpha_6(1800),N_6(1800))=(582,1152),
\]
in agreement with Proposition~\ref{prop:same-support}.  These exact
counts are finite illustrations of the proved path-cover mechanism, not
asymptotic evidence.
\end{remark}

The finite MWIS and square-clock viewpoint follows the combinatorial
template of RH-375, Theorem~2.2 and Sections~3--5
\cite{RH375}.  The \(h=1\) centered marginal argument in RH-395 is the
direct finite precedent \cite{RH395}.  Neither predecessor supplies the
general distance-\(d\), collision-aware proof above, and neither is used
as a terminal-clock analytic input.

\section{Euler--run formula}

For a finite prime set \(\mathcal P\) and finite \(J\subset\Z\), define
\begin{equation}
 D_{h,\mathcal P}(J)=
 \prod_{p\in\mathcal P}\left(
 1-\frac{|\{dj\bmod p^2:j\in J\}|}{p^2}\right),
 \label{eq:finite-D}
\end{equation}
and define \(R_{\ell,h,\mathcal P}\) by the same four-term expression as
\eqref{eq:Rlh}, with \(D_h\) replaced by
\(D_{h,\mathcal P}\).  Thus
\begin{align}
 R_{\ell,h,\mathcal P}
 &=D_{h,\mathcal P}([0,\ell-1])
 -D_{h,\mathcal P}(\{-1\}\cup[0,\ell-1])\notag\\
 &\quad-D_{h,\mathcal P}([0,\ell])
 +D_{h,\mathcal P}([-1,\ell]).
 \label{eq:finite-R}
\end{align}

\begin{lemma}[Exact finite run densities]
\label{lem:finite-runs}
Suppose \(p_0(h)\in\mathcal P\).  On
\(\Z/q_{\mathcal P}\Z\), the exact number of bracketed positive
step-\(d\) runs of length \(\ell\) is
\begin{equation}
 q_{\mathcal P}R_{\ell,h,\mathcal P}.
 \label{eq:run-count}
\end{equation}
These quantities are nonnegative integers and vanish for
\(\ell\geq p_0(h)^2\).  Moreover,
\begin{equation}
 \alpha_h(q_{\mathcal P})
 =\frac{N_h(q_{\mathcal P})}{2}
 +\frac{q_{\mathcal P}}2
 \sum_{\substack{1\leq\ell<p_0(h)^2\\\ell\ {\rm odd}}}
 R_{\ell,h,\mathcal P}.
 \label{eq:alpha-run}
\end{equation}
\end{lemma}

\begin{proof}
For \(r\bmod q_{\mathcal P}\), all phases
\(r+dj\), \(j\in J\), are positive precisely when, for each
\(p\in\mathcal P\), the residue \(r\pmod{p^2}\) avoids the distinct
classes \(-dj\pmod{p^2}\).  The Chinese remainder theorem therefore gives
the exact density \eqref{eq:finite-D}.  Inclusion--exclusion on the two
endpoint events \(r-d\) positive and \(r+d\ell\) positive gives
\eqref{eq:finite-R}.  Its complement interpretation is exactly: zero at
index \(-1\), positive at indices \(0,\ldots,\ell-1\), and zero at index
\(\ell\).  This proves \eqref{eq:run-count} and nonnegativity.

Because \(p_0\nmid d\), the residues \(dj\pmod{p_0^2}\) for any
\(p_0^2\) consecutive indices \(j\) cover all classes.  Every such block
therefore contains a phase divisible by \(p_0^2\), proving the run cutoff.
There are no all-positive cycles.

Finally, a path of length \(\ell\) has MWIS cardinality
\[
 \left\lceil\frac\ell2\right\rceil
 =\frac\ell2+\frac12\ind_{\{\ell\ {\rm odd}\}}.
\]
Summing the first term over all runs gives \(N_h/2\), and summing the
second counts the odd runs.  Equation \eqref{eq:run-count} now gives
\eqref{eq:alpha-run}.
\end{proof}

\begin{proposition}[Finite Euler--run endpoint]
\label{prop:finite-endpoint}
If \(p_0(h)\in\mathcal P\), then
\begin{align}
 C_h(q_{\mathcal P})
 &=M_h(q_{\mathcal P})=B_{h,\mathcal P},
 \label{eq:finite-B-equality}\\
 B_{h,\mathcal P}
 &=\frac{K_1}{2}
 +\frac{K_1}{2D_{h,\mathcal P}(\{0\})}
 \sum_{\substack{1\leq\ell<p_0(h)^2\\\ell\ {\rm odd}}}
 R_{\ell,h,\mathcal P}.
 \label{eq:finite-B}
\end{align}
\end{proposition}

\begin{proof}
The singleton finite density satisfies
\[
 D_{h,\mathcal P}(\{0\})
 =\prod_{p\in\mathcal P}\left(1-\frac1{p^2}\right)
 =\frac{N_h(q_{\mathcal P})}{q_{\mathcal P}}.
\]
Insert \eqref{eq:alpha-run} into the normalized definition
\eqref{eq:MQ}.  This gives \eqref{eq:finite-B}.
Proposition~\ref{prop:square-saturation} gives
\eqref{eq:finite-B-equality}.
\end{proof}

\begin{lemma}[Infinite densities and termwise limit]
\label{lem:infinite-density}
For every fixed finite \(J\subset\Z\), the products
\(D_{h,\mathcal P}(J)\) converge to \(D_h(J)\) as
\(\mathcal P\) increases through the primes.  The limit is the density of
simultaneous squarefreeness at the step-\(d\) positions indexed by \(J\).
Consequently
\[
 R_{\ell,h,\mathcal P}\longrightarrow R_{\ell,h}\geq0
 \qquad(\ell\text{ fixed}).
\]
\end{lemma}

\begin{proof}
The number of excluded classes at a prime is at most \(|J|\), and
\[
 \sum_p\frac{|\{dj\bmod p^2:j\in J\}|}{p^2}
 \leq |J|\sum_p\frac1{p^2}<\infty.
\]
Thus the local products converge absolutely in the standard product
sense (with a possible zero caused only by a finite local factor).
Finite CRT gives every truncated density.  Removing the truncation is
justified by the square-divisor union bound
\[
 \overline d\{n:p^2\mid n+dj
 \text{ for some }j\in J,\ p>Y\}
 \leq |J|\sum_{p>Y}\frac1{p^2}\longrightarrow0.
\]
Hence the product is the claimed density.  Applying the same
two-endpoint inclusion--exclusion as in
Lemma~\ref{lem:finite-runs} shows that the four-term limit is an exact
bracketed-run density, so it is nonnegative.
\end{proof}

\begin{proposition}[Cofinal convergence]
\label{prop:cofinal}
Let \(\mathcal P_y\) be any increasing prime-initial family containing
\(p_0(h)\) and tending to all primes.  Then
\begin{equation}
 B_{h,\mathcal P_y}\longrightarrow B_\infty(h).
 \label{eq:cofinal-limit}
\end{equation}
\end{proposition}

\begin{proof}
Lemma~\ref{lem:infinite-density} gives
\[
 D_{h,\mathcal P_y}(\{0\})\longrightarrow
 D_h(\{0\})=\prod_p(1-p^{-2})=K_1.
\]
It also gives convergence of every run-density term.  The run cutoff
\(\ell<p_0(h)^2\) is fixed and finite, so the limit can be taken
termwise in \eqref{eq:finite-B}.  Its first term is
\(K_1/2=3/\pi^2\), and its run coefficient tends to \(1/2\).  The result
is exactly \eqref{eq:Binfinity}.
\end{proof}

\begin{remark}[Certified numerical orientation only]
The exact rational interval oracle in the companion artifact encloses
\[
 B_\infty(1)\in[0.4214,0.4224],\qquad
 B_\infty(2)\in[0.3282,0.3296],\qquad
 B_\infty(3)\in[0.4143,0.4181].
\]
These intervals orient the scale of the endpoint.  They are not used in
any comparison or proof, and they imply no monotonicity in \(h\).
\end{remark}

\section{Fresh-prime lifts and strict nonattainment}

Fix a finite prime support \(\mathcal P\) containing \(p_0(h)\), and let
\(\varpi\notin\mathcal P\) be a prime with
\(\varpi\nmid d q_{\mathcal P}\).  Put
\(\mathcal P'=\mathcal P\cup\{\varpi\}\).  Every old positive run has
length less than \(p_0^2\), and
\[
 p_0^2<\varpi^2,
\]
because \(p_0\) is already supported and no smaller prime can be both
fresh and coprime to \(d\).

\begin{lemma}[Deletion parity on one lifted path]
\label{lem:path-deletion}
Let \(P_L\) be a path with vertices \(1,\ldots,L\).  Deleting vertex
\(j\) leaves MWIS cardinality
\begin{equation}
 \left\lceil\frac{j-1}{2}\right\rceil
 +\left\lceil\frac{L-j}{2}\right\rceil.
 \label{eq:deleted-path}
\end{equation}
If \(L\) is odd, this is one less than
\(\lceil L/2\rceil\) exactly when \(j\) is odd.  If \(L\) is even, it
always equals \(L/2\).
\end{lemma}

\begin{proof}
Deleting \(j\) splits \(P_L\) into paths of lengths \(j-1\) and \(L-j\);
their MWIS cardinalities add, giving \eqref{eq:deleted-path}.  Substitution
of \(L=2m+1\) or \(L=2m\) gives the two parity statements.
\end{proof}

\begin{proposition}[Fresh-prime recurrence]
\label{prop:fresh-recurrence}
Let \(\mathcal R_{\mathcal P}\) be the multiset of old positive run
lengths and set
\[
 E_{\mathcal P}=
 \sum_{\substack{L\in\mathcal R_{\mathcal P}\\L\ {\rm even}}}\frac L2,
 \qquad
 O_{\mathcal P}=
 \sum_{\substack{L\in\mathcal R_{\mathcal P}\\L\ {\rm odd}}}
 \left\lceil\frac L2\right\rceil.
\]
Then
\begin{align}
 N_h(q_{\mathcal P'})
 &=(\varpi^2-1)N_h(q_{\mathcal P}),\label{eq:Nprime}\\
 \alpha_h(q_{\mathcal P'})
 &=\varpi^2\alpha_h(q_{\mathcal P})-O_{\mathcal P}\notag\\
 &=(\varpi^2-1)\alpha_h(q_{\mathcal P})+E_{\mathcal P},
 \label{eq:alphaprime}\\
 B_{h,\mathcal P'}-B_{h,\mathcal P}
 &=\frac{K_1E_{\mathcal P}}
 {(\varpi^2-1)N_h(q_{\mathcal P})}.
 \label{eq:normalized-gain}
\end{align}
Thus a fresh-prime step is strict exactly when the old positive graph has
an even run.  It is not necessarily strict at every step.
\end{proposition}

\begin{proof}
Because \(\varpi\nmid q_{\mathcal P}\), every old phase has
\(\varpi^2\) lifts modulo \(q_{\mathcal P'}\).  Exactly one lift is
divisible by \(\varpi^2\), proving \eqref{eq:Nprime}.

Consider an old run of length \(L<\varpi^2\).  Among its
\(\varpi^2\) lifted copies, \(\varpi^2-L\) have no deletion, and for
each \(j=1,\ldots,L\) exactly one copy has its \(j\)th vertex deleted.
Indeed, the lift parameter runs bijectively modulo \(\varpi^2\); two
different positions cannot be deleted in the same lift because that
would give
\(\varpi^2\mid d(j_1-j_2)\), while
\(\varpi\nmid d\) and \(0<|j_1-j_2|<\varpi^2\).

Lemma~\ref{lem:path-deletion} shows that an odd run loses one unit for
each odd deletion position, a total loss \(\lceil L/2\rceil\), whereas
an even run loses nothing.  Its total lifted contribution is therefore
\[
 \begin{cases}
  (\varpi^2-1)\lceil L/2\rceil,&L\text{ odd},\\
  \varpi^2L/2,&L\text{ even}.
 \end{cases}
\]
Summing runs gives both forms of \eqref{eq:alphaprime}.  Dividing
\eqref{eq:alphaprime} by \eqref{eq:Nprime} in
\eqref{eq:MQ} proves \eqref{eq:normalized-gain}.
\end{proof}

\begin{remark}[A genuine plateau]
\label{rem:plateau}
For \(h=9\), exact raw counts give
\[
 (\alpha_9(36),N_9(36))=(16,24),\qquad
 (\alpha_9(900),N_9(900))=(384,576),
\]
and therefore
\[
 M_9(36)=M_9(900)=\frac{2K_1}{3}.
\]
The first clock does not yet contain \(p_0(9)^2=25\), so this fixture is
outside the recurrence's qualified base domain.  It nevertheless rules
out any blanket claim of strict gain at every square-support prime step.
Proposition~\ref{prop:fresh-recurrence}, rather than stepwise strictness,
is the exact statement.
\end{remark}

\begin{lemma}[CRT creation of an even run]
\label{lem:CRT-two-run}
Every finite support containing \(p_0(h)\) has a finite extension whose
positive graph contains an exact run of length two.
\end{lemma}

\begin{proof}
Enlarge the support, if necessary, until it contains two distinct primes
\(a,b\) with \(a\nmid d\) and \(b\nmid d\).  Seek a starting phase \(r\)
whose interior phases \(r,r+d\) are positive and whose endpoints
\(r-d,r+2d\) are zero.  Impose
\[
 r\equiv d\pmod{a^2},\qquad
 r\equiv-2d\pmod{b^2}.
\]
The left endpoint is then divisible by \(a^2\), and the right endpoint
is divisible by \(b^2\).  Since neither prime divides \(d\), the two
interior phases are nonzero modulo both prime squares.  At every other
supported prime \(p\), choose a local residue avoiding the two classes
\(0,-d\pmod{p^2}\); at least two of the \(p^2\) classes remain even for
\(p=2\).  The Chinese remainder theorem combines these choices.  The
resulting \(r\) begins an exact positive run of length two.
\end{proof}

\begin{proposition}[Eventual strictness]
\label{prop:eventual-strictness}
For every finite \(\mathcal P\) containing \(p_0(h)\), some finite
extension \(\mathcal P^\star\) satisfies
\[
 B_{h,\mathcal P^\star}>B_{h,\mathcal P}.
\]
Consequently \(B_{h,\mathcal P}<B_\infty(h)\) whenever
\(\mathcal P\) also contains every prime divisor of \(d\).
\end{proposition}

\begin{proof}
Add generic primes until Lemma~\ref{lem:CRT-two-run} applies.  Every such
addition is nondecreasing by \eqref{eq:normalized-gain}.  The extended
graph now has an even run, so adjoining one further fresh prime is
strict, again by \eqref{eq:normalized-gain}.  This proves the first
assertion.

For the second, extend \(\mathcal P\) through the missing primes in
increasing order.  Since all prime divisors of \(d\) were already
included, every newly added prime is coprime to \(d\), so the recurrence
applies at every step.  The values are nondecreasing, some future step is
strict, and Proposition~\ref{prop:cofinal} identifies their limit as
\(B_\infty(h)\).  Hence the starting value is strictly below that limit.
\end{proof}

\begin{proof}[Proof of Theorem~\ref{thm:endpoint}]
First let \(\mathcal P_y\) be a cofinal prime-initial family containing
\(p_0(h)\).  Proposition~\ref{prop:finite-endpoint} and
\eqref{eq:cofinal-limit} give
\[
 C_h(q_{\mathcal P_y})=B_{h,\mathcal P_y}
 \longrightarrow B_\infty(h).
\]
Therefore \(\sup_q C_h(q)\geq B_\infty(h)\).

For the strict upper bound, fix an arbitrary finite \(q\).  Choose a
finite prime set \(\mathcal P\) containing \(p_0(h)\), every prime divisor
of \(q\), and every prime divisor of \(d\).  Set
\[
 Q=\operatorname{lcm}(q,q_{\mathcal P}).
\]
Then \(q\mid Q\), \(q_{\mathcal P}\mid Q\), and
\(\rad(Q)=\rad(q_{\mathcal P})\).  Lemma~\ref{lem:clock-divisibility},
Proposition~\ref{prop:same-support}, and
Proposition~\ref{prop:eventual-strictness} give
\[
 C_h(q)\leq C_h(Q)=M_h(q_{\mathcal P})
 =B_{h,\mathcal P}<B_\infty(h).
\]
This proves strict nonattainment at every finite clock and the reverse
supremum inequality.  At no point is \(q\) allowed to depend on \(X\);
the scalar supremum is taken only after all fixed-table limits and finite
maxima.
\end{proof}

\section{The landscape across fixed lags}

\begin{lemma}[An isolated run has positive density]
\label{lem:isolated-positive}
For every fixed \(h\geq1\),
\[
 R_{1,h}>0.
\]
\end{lemma}

\begin{proof}
Choose distinct primes \(a,b\) with \(a\nmid d\) and \(b\nmid d\).
The Chinese remainder theorem gives one residue class modulo
\(a^2b^2\) satisfying
\begin{equation}
 r\equiv d\pmod{a^2},\qquad
 r\equiv-d\pmod{b^2}.
 \label{eq:isolated-congruences}
\end{equation}
Every integer in this class has \(a^2\mid r-d\) and
\(b^2\mid r+d\).  Moreover, its center is nonzero modulo both
\(a^2\) and \(b^2\), because neither prime divides \(d\).  Among these
integers, impose for each remaining prime \(p\) the local avoidance
\(r\not\equiv0\pmod{p^2}\).  Finite CRT counts followed by the
elementary square-divisor tail give this subevent the density
\[
 \frac1{a^2b^2}
 \prod_{p\ne a,b}\left(1-\frac1{p^2}\right)>0.
\]
Here the product is positive because \(\sum_p p^{-2}<\infty\).  On the
subevent, the center \(r\) is squarefree while the two step-\(d\)
neighbors \(r-d\) and \(r+d\) are not.  It is therefore contained in
the exact bracketed length-one event whose density is \(R_{1,h}\).
\end{proof}

\begin{lemma}[Outside-prime boundary bound]
\label{lem:outside-prime-tail}
For every integer \(Y\geq2\), put
\begin{equation}
 d_Y=\prod_{\substack{p\leq Y\\p\ {\rm prime}}}p^2,
 \qquad h_Y=\frac{d_Y}{2}.
 \label{eq:hY}
\end{equation}
Then \(h_Y\) is a positive integer and
\begin{equation}
 0\leq B_\infty(h_Y)-\frac3{\pi^2}
 \leq\frac12\sum_{p>Y}\frac1{p^2}
 <\frac1{2(Y-1)}.
 \label{eq:outside-tail}
\end{equation}
\end{lemma}

\begin{proof}
The factor \(2^2\) occurs in \(d_Y\), so \(h_Y\) is integral and
\(2h_Y=d_Y\).  Exact positive runs in step-\(d_Y\) coordinates
partition the run-start event
\[
 \{r:\ r\text{ is squarefree and }r-d_Y\text{ is not squarefree}\}.
\]
All runs are finite: the first prime not dividing \(d_Y\) supplies the
run-length cutoff proved above.  Consequently the density of this
event is
\begin{equation}
 \sum_{\ell\geq1}R_{\ell,h_Y}.
 \label{eq:all-run-starts}
\end{equation}
If \(p\leq Y\) and \(p^2\mid r-d_Y\), then \(p^2\mid d_Y\) forces
\(p^2\mid r\), contrary to the center being squarefree.  Every boundary
in \eqref{eq:all-run-starts} must therefore be supplied by a prime
\(p>Y\).  The union bound and the density \(p^{-2}\) of one square
divisibility class give
\[
 \sum_{\ell\geq1}R_{\ell,h_Y}
 \leq\sum_{p>Y}\frac1{p^2}.
\]
The endpoint bonus uses only odd runs, so nonnegativity of all the
\(R_{\ell,h_Y}\) and \eqref{eq:Binfinity} prove the first inequality in
\eqref{eq:outside-tail}.  Finally,
\(\sum_{p>Y}p^{-2}\leq\sum_{n>Y}n^{-2}<1/(Y-1)\), which proves the
displayed strict bound.
\end{proof}

\begin{proof}[Proof of Corollary~\ref{cor:lag-infimum}]
All run densities in \eqref{eq:Binfinity} are nonnegative, so
\(B_\infty(h)\geq3/\pi^2\).  Lemma~\ref{lem:isolated-positive} puts the
strictly positive term \(R_{1,h}/2\) into the endpoint for every fixed
\(h\), because \(1<p_0(h)^2\).  Hence equality is impossible at any
fixed lag.

On the other hand, Lemma~\ref{lem:outside-prime-tail} supplies a
sequence of individually fixed integer lags \(h_Y\) for which the
endpoint bonus tends to zero.  Thus values approach \(3/\pi^2\) from
above, proving \eqref{eq:lag-infimum} and nonattainment.  This argument
compares scalar endpoints only after their fixed-lag limits; it provides
neither a varying-lag terminal limit nor an ordering of endpoints by
lag.
\end{proof}

\section{Source roles, reproducibility, and limitations}

\subsection{Analytic and combinatorial provenance}

The sole analytic input is the complete three-shift terminal-log law of
RH-394, frozen at repository commit
\nolinkurl{6b3d616851cd2d7cba66371d0aa9f25b8e8bf2f7}
\cite{RH394}.  Section~2 uses precisely its fixed-shift density,
phase-average, mass, and terminal-clock conclusions.  The present paper
supplies every subsequent safety, tropical, square-clock, run-density,
strictness, and lag-landscape argument.

RH-395, frozen at commit
\nolinkurl{20de7202518f4488cbd9c7d63bf94aaa3dc94476}, is cited only as
the \(h=1\) finite relation-saturation and tropical-optimizer precedent
\cite{RH395}.  RH-375, frozen at commit
\nolinkurl{071fed1b2a5d8488b9d2e35a99a753953b233584}, is cited only as
the finite one-site MWIS and square-clock template \cite{RH375}.
Neither provides an analytic terminal-clock input here.  The
Tao two-point and Tao--Ter\"av\"ainen odd-correlation papers are
inherited provenance through RH-394 \cite{Tao2016LogChowla,
TaoTeravainen2019}; the Johnston--Yang and Maynard records likewise
remain inherited source-closure objects only.  No new conclusion in
this paper is attributed directly to those four remote records.

\subsection{Finite reproduction and source integrity}

The companion certificate is an exact finite reproduction artifact,
not an analytic proof.  Its frozen canonical JSON has \(96\) rows,
\(83{,}309\) bytes, and SHA-256
\nolinkurl{7cc0da78ee7e47a22b357d7e8d907bc9d9879caeb82ede30709e8cb1023032ba}.
The strict validator rejects all \(32\) named single-core mutations.
Among its checks, the exhaustive full-eight relation oracle scans
\(262{,}144\) ordered relation pairs and finds exactly \(3{,}375\)
universally safe pairs.  These checks reproduce finite identities and
counterexamples used above; they do not replace the proofs.

The frozen source closure contains \(160\) immutable Git objects and
four inherited remote records, hence \(164\) logical objects.  Its
canonical SHA-256 is
\nolinkurl{c16456d58efd74edf1505c430a54459e359b5ba7e1e581773e9a0613b493385b};
the all-Git and logical-source digests are, respectively,
\nolinkurl{472bf5ce5e352dce0d3a44ad10b22345b98e0e8b9a0cd745be9ecd93dedf0a86}
and
\nolinkurl{72040ab3d7a5d98ce308b91d0748d52a8d4886cf245f5079f14c69ee659cc287}.
Reproduction is offline: the release records zero network requests,
vendors no external PDF payload, and records no forbidden payload-hash
hit.  The four remote redistribution flags are
\((\mathrm{no},\mathrm{no},\mathrm{yes},\mathrm{no})\); the corresponding
PDFs are not redistributed.

\subsection{Claim firewall}

Every analytic assertion fixes \(h\), \(q\), and the complete phase
table before \(X\to\infty\), and it holds for every terminal clock in
\eqref{eq:clock}.  The order is: terminal limit for each fixed table,
finite safe maximum, and only then the scalar supremum over finite
clocks.  The result is not an ordinary Ces\`aro-average theorem and
contains no regime with \(h=h(X)\), \(q=q(X)\), or growing phase tables;
no uniform rate in these parameters is asserted.

The centered rule reads \(\mob(n+h)\), so it is not causal or online.
The proof concerns exactly three shifts and safety at the single
separation \(2h\); it does not establish an even four-shift law, a
larger-window compiler, or a generic graph theorem.  Nothing here is an
analytic trace or operator formula, a statement about zeta zeros or the
Riemann hypothesis, or a discharge of downstream Gates A--E.  Finally,
the lag argument proves only the infimum and its nonattainment.  It
makes no assertion about a supremum, a maximum, or monotonicity in
\(h\).

\section*{Declarations}

\noindent\textbf{Data availability.}
The finite certificate, exact result and schema, source-lock metadata, and
offline verification code are included in the accompanying RH-396 package.
No new empirical data were generated and no external source PDF is
redistributed.

\noindent\textbf{Ethics declaration.}
This mathematical study involved no human participants, animals, personal
data, or sensitive datasets; institutional ethics approval was not
applicable.

\noindent\textbf{Author contributions (CRediT).}
The RH research program performed conceptualization, formal analysis,
methodology, software, validation, and writing (original draft, review,
and editing).

\noindent\textbf{Conflict of interest.}
The author declares no conflict of interest.

\noindent\textbf{Funding.}
No external funding was received for this work.

\noindent\textbf{AI-use statement.}
AI-assisted tools supported drafting, finite-certificate implementation,
formatting, and consistency checks.  The research program directed and
reviewed the mathematical arguments, source roles, and release artifacts
and retains responsibility for the content.

\begingroup
\sloppy
\raggedright
\bibliographystyle{plainnat}
\bibliography{references}
\endgroup

\end{document}
